% Einheit 3 von 5 – Dreieck (bank/flaechen/e3.jsonl, zusammenbau v0.1)
%% TODO Zweigzeile Teil 1: was hier gelernt wird (Satz in Schülersprache aus dem Einheitstitel) – Einheit 3
\einheitenkopf[e3]{Einheit 3 von 5 · Dreieck}
\zweigzeile{ab Kl. 7, je nach Buch bis Kl. 8 (am Gymnasium ab Kl. 5) · P10 oft}
\setcounter{aufgabe}{17}

%% TODO Titel: Ich-kann-Satz fehlt in der Bank (Platzhalter: Kettenname) – Nr. 18
%% TODO Anweisung: ein Satz über der ersten Teilaufgabe fehlt in der Bank – Nr. 18
\begin{aufgabe}{Dreieck}
\swa{Zeichne die Höhe zur Grundseite $g$ ein und markiere den rechten Winkel.}{\dreieck{(0,0)}{(5,0)}{(2,3)}{}{}{$g$}{}{}{}}
%% TODO Darstellung für schwach fehlt in der Bank (flaechen-e3-k1-s1-v1; Abschnitt „Für schwache Schüler“)
\swz{Dreieck, $g = 8$ cm, $h = 7$ cm – Fläche?}{}
%% TODO Darstellung für schwach fehlt in der Bank (flaechen-e3-k1-s1-v2; Abschnitt „Für schwache Schüler“)
\swz{Dreieck, $g = 14$ m, $h = 9$ m – Fläche?}{}
%% TODO Darstellung für schwach fehlt in der Bank (flaechen-e3-k1-s1-v3; Abschnitt „Für schwache Schüler“)
\swz{Ein dreieckiges Tuch: $g = 6$ dm, $h = 11$ dm – Fläche?}{}
%% TODO Darstellung für schwach fehlt in der Bank (flaechen-e3-k1-s1-v4; Abschnitt „Für schwache Schüler“)
\swz{Dreieck, $g = 20$ cm, $h = 15$ cm – Fläche?}{}
%% TODO Darstellung für schwach fehlt in der Bank (flaechen-e3-k1-s1-v5; Abschnitt „Für schwache Schüler“)
\swz{Ein dreieckiges Beet: $g = 9$ m, $h = 5$ m – Fläche?}{}
\swfrage{Was bleibt gleich, was ändert sich?}
\swz{Rechtwinkliges Dreieck – Fläche?}{\dreieckrw{5}{2.67}{$15$ cm}{$8$ cm}{$17$ cm}}
\swz{Stumpfwinkliges Dreieck, Grundseite $6$ cm; die Höhe dazu liegt außerhalb und ist $4$ cm lang – Fläche?}{\dreieck{(0,0)}{(3,0)}{(4.5,2)}{}{}{$6$ cm}{}{}{}}
%% TODO Darstellung für schwach fehlt in der Bank (flaechen-e3-k1-s4-v1; Abschnitt „Für schwache Schüler“)
\swz{Dreieck, $g = 4{,}6$ cm, $h = 3$ cm – Fläche?}{}
%% TODO Darstellung für schwach fehlt in der Bank (flaechen-e3-k1-s5-v1; Abschnitt „Für schwache Schüler“)
\swz{Dreieck, $A = 26$ cm², $h = 4$ cm – Grundseite?}{}
%% TODO Darstellung für schwach fehlt in der Bank (flaechen-e3-k1-s6-v1; Abschnitt „Für schwache Schüler“)
\swz{Ein gleichseitiges Dreieck hat die Seitenlänge $s$. Term für den Umfang?}{}
\swz[3]{Ein dreieckiges Grundstück ABC hat bei C einen rechten Winkel; $BC = 16$ m, $AC = 45$ m. Ein Teilstück QBC mit Q auf AC soll $120$ m² groß sein. Wie lang ist QC? \hfill (P10 2020 OS)}{}
\end{aufgabe}

%% TODO Titel: Ich-kann-Satz fehlt in der Bank (Platzhalter: Kettenname) – Nr. 19
%% TODO Anweisung: ein Satz über der ersten Teilaufgabe fehlt in der Bank – Nr. 19
\begin{aufgabe}{Höhe zur passenden Grundseite wählen (drei Höhen)}
%% TODO Darstellung für schwach fehlt in der Bank (flaechen-e3-k2-s1-v1; Abschnitt „Für schwache Schüler“)
\swz{Dreieck ABC: $AB = 8$ cm, $BC = 7{,}5$ cm. Die Höhe auf AB ist $6$ cm, die Höhe auf BC $6{,}4$ cm. Fläche – rechne mit BC und der passenden Höhe.}{}
\end{aufgabe}

%% TODO Titel: Ich-kann-Satz fehlt in der Bank (Platzhalter: Kettenname) – Nr. 20
%% TODO Anweisung: ein Satz über der ersten Teilaufgabe fehlt in der Bank – Nr. 20
\begin{aufgabe}{Umfang}
%% TODO Darstellung für schwach fehlt in der Bank (flaechen-e3-k3-s1-v1; Abschnitt „Für schwache Schüler“)
\swz{Dreieck mit den Seiten $5$ cm, $7$ cm und $9$ cm – Umfang?}{}
\end{aufgabe}

%% TODO Titel: Ich-kann-Satz fehlt in der Bank (Platzhalter: Kettenname) – Nr. 21
%% TODO Anweisung: ein Satz über der ersten Teilaufgabe fehlt in der Bank – Nr. 21
\begin{aufgabe}{Dreieck – Fehler finden, Begründen, Anwendung, Darstellungswechsel}
\swz[3]{Dreieck mit $g = 9$ cm und $h = 6$ cm. Jana rechnet: \rechnung{A &= 9 \cdot 6 \\ A &= 54 \text{ cm}^2} Wo ist der Fehler?}{}
\swfrage{Warum rechnet man beim Dreieck Grundseite mal Höhe und nimmt davon die Hälfte? Begründe mit einem Parallelogramm.}
\swz[3]{Ein dreieckiger Hausgiebel ist unten $8{,}4$ m breit und $3{,}5$ m hoch. Er wird gestrichen; eine Dose Farbe reicht für $6$ m². Wie viele Dosen braucht man?}{}
\swa{Zeichne ein Dreieck, dessen Umfang $u = 3 \cdot k$ ist, und beschrifte die Seiten.}{\rechenplatz[halb]{4}}
\end{aufgabe}

\uebersichtskasten{Fläche: Grundseite mal Höhe, davon die Hälfte. Rechtwinklig: die Katheten sind Grundseite und Höhe. \\ g = 10 cm, h = 6 cm: A = 10 · 6 : 2 = 30 cm² \quad Katheten 3 cm und 4 cm: A = 3 · 4 : 2 = 6 cm² \\ Rückwärts: g = 2 · A : h. \quad A = 30 cm², h = 6 cm: g = 60 : 6 = 10 cm}
